% =============================================================================
%  IEEE ISAC 2026 — Conference Paper Draft
%  Title: Data-Driven Dynamical Models for Rainfall Nowcasting from
%         Commercial Microwave Links
%
%  Build instructions (Overleaf):
%    1. New Project -> Upload Project -> drop this .tex in (or paste contents).
%    2. Compiler: pdfLaTeX. Bibliography is inline (\bibitem) -- no .bib needed.
%    3. Results numbers are marked \tbd{} -- replace once experiments finish.
% =============================================================================

\documentclass[conference]{IEEEtran}
\IEEEoverridecommandlockouts

\usepackage{cite}
\usepackage{amsmath,amssymb,amsfonts}
\usepackage{algorithmic}
\usepackage{textcomp}
\usepackage{xcolor}
\usepackage{booktabs}
\usepackage{array}
\usepackage{multirow}
\usepackage{url}
\usepackage{hyperref}

\newcommand{\tbd}{\textcolor{red}{\textbf{TBD}}}
\newcommand{\note}[1]{\textcolor{blue}{[#1]}}

\def\BibTeX{{\rm B\kern-.05em{\sc i\kern-.025em b}\kern-.08em
    T\kern-.1667em\lower.7ex\hbox{E}\kern-.125emX}}

\begin{document}

% =============================================================================
\title{Data-Driven Dynamical Models for Rainfall Nowcasting from Commercial Microwave Links}

\author{%
\IEEEauthorblockN{Author Name\IEEEauthorrefmark{1}}
\IEEEauthorblockA{%
\IEEEauthorrefmark{1}Department / Lab\\
Institution\\
City, Country\\
\texttt{email@domain}}%
}

\maketitle

% =============================================================================
\begin{abstract}
Commercial microwave link (CML) networks deployed by cellular operators
form a natural \emph{opportunistic} Integrated Sensing and Communication
(OISAC) substrate: links already provisioned for backhaul attenuate
during rainfall and can be repurposed as a dense rain-sensing layer at
zero hardware cost, without any modification of the underlying
communication waveform. While CML-based
rainfall \emph{estimation} is by now a mature line of work, short-term
rainfall \emph{nowcasting} (15--60\,min ahead) from CML observations is
an emerging direction in which a unifying methodological framework that
places competing dynamical paradigms on equal footing is still missing.
We propose such a framework as a two-stage spatiotemporal learning
problem: a recurrent estimation network maps raw CML attenuation signals
to per-link rain rates which are interpolated onto a regular grid, and a
downstream \emph{dynamical model} learns to propagate the resulting
CML-derived rain fields forward in time. A central methodological
proposition is a \emph{CML self-supervised} forecasting objective in
which the dynamical model is trained against future CML-derived fields
with no external rain reference at Stage~2, preserving the OISAC
premise end to end. Within this framework we benchmark three families of data-driven
dynamical models against persistence baselines: a Transformer encoder
over flattened map tokens, a reduced-order POD--SINDy model that
identifies an explicit nonlinear ODE in a low-dimensional latent space,
and a recurrent state-space sequence model in the spirit of Mamba. Each
family is further ablated under single-step versus multi-horizon
prediction heads and under the proposed CML self-supervision versus a
sparse gauge-supervised baseline. Experiments
on the OpenMRG dataset (Gothenburg, summer 2015) with both point-to-pixel
(against rain gauges) and grid-to-grid (against weather radar) evaluation
indicate that \tbd{}.
\end{abstract}

\begin{IEEEkeywords}
Integrated sensing and communication, opportunistic sensing, commercial
microwave links, rainfall nowcasting, spatiotemporal forecasting,
state-space models, transformers, sparse identification of nonlinear
dynamics.
\end{IEEEkeywords}

% =============================================================================
\section{Introduction}
\label{sec:intro}
While ISAC as a co-design paradigm targets future beyond-5G systems, opportunistic weather sensing with existing wireless communication network (WCN) infrastructure is already an established, operational reality. The foundational concept of repurposing communication signals for environmental monitoring was initially demonstrated using commercial microwave links (CMLs)~\cite{messer2006}. Operating in the $6-80$ GHz band and widely deployed as cellular backhaul, these CMLs serve as a canonical ISAC modality. Their received signal level (RSL) is directly attenuated by hydrometeors along the path, providing a highly informative measure of path-averaged rain rates~\cite{leijnse2007,overeem2013}. This approach now leverages infrastructure across diverse network types and geographical scales. Regional and city-scale datasets, such as OpenRainER (Emilia-Romagna, Italy)~\cite{covi2026openrainer} and OpenMRG (Gothenburg, Sweden)~\cite{andersson2022openmrg}, provide detailed coverage over varied terrain. Meanwhile, country-wide efforts include multi-year CML archives for the Netherlands~\cite{overeem2024netherlands} and operational rainfall monitoring across Germany~\cite{graf2020rainfall}, which are typically verified using radar products or official weather stations. Moreover, as next-generation networks increasingly deploy high-frequency, short-range links in V-band (60 GHz) and E-band (70–80 GHz) to achieve dense urban coverage, they unlock unprecedented spatial resolution for environmental monitoring. The utilization of community-driven mesh networks, such as those in New York City~\cite{jacoby2025openmesh}, alongside smart-city infrastructure~\cite{hadar2020parameter,ostrometzky2024opportunistic}, demonstrates high-resolution urban weather sensing as a rapidly emerging application of this ubiquitous architecture. WCN therefore constitute a free, dense, and distributed OISAC weather-sensing layer that complements traditional monitoring systems and has been recognized as a key component in ISAC frameworks for environmental monitoring~\cite{messer2006,cui2023integrated}.
% Therefore, WCN backhaul networks, which operating at microwave and mmWave frequencies, constitute a free, dense, and distributed OISAC weather-sensing layer that complements traditional monitoring systems, with growing recognition in ISAC research for environmental sensing applications~\cite{9705498,cui2023integrated}.

The bulk of CML-related sensing research has so far focused on the
\emph{estimation} problem: given CML attenuation at time $t$, recover
the instantaneous rain rate at $t$~\cite{habi2020,zhang2023environment,chwala2019}. The
\emph{nowcasting} problem -- predicting rainfall fields several tens of
minutes into the future from CML observations alone -- has received much
less attention, despite being arguably more useful for downstream
applications such as urban drainage, flash-flood early warning, and
short-horizon decision support for outdoor activity. Weather radar, while accurate, requires significant capital investment and operational costs, and provides limited coverage in many regions. Rain gauges offer direct point measurements but are extremely sparse and fail to capture spatial variability.

Modern data-driven sequence modelling offers three broadly distinct
families of dynamical models that are natural candidates for this task.
\emph{Attention-based} models such as the
Transformer~\cite{vaswani2017} learn unconstrained long-range
dependencies among input tokens. \emph{State-space sequence models},
including Mamba~\cite{gu2023mamba}, parametrise a continuous-time linear
dynamical system whose parameters depend on the input, yielding
recurrent inference with strong long-context behaviour. \emph{Equation
discovery} methods such as the Sparse Identification of Nonlinear
Dynamics (SINDy)~\cite{brunton2016sindy}, combined with a Proper
Orthogonal Decomposition (POD), instead identify an \emph{explicit}
parsimonious ODE in a low-dimensional latent space. Each family encodes
a different inductive bias about how a rainfall field evolves; a fair
head-to-head comparison on a CML-driven nowcasting task is, to our
knowledge, missing from the literature.

\paragraph*{Contributions} This paper makes four contributions:

\begin{itemize}
  \item We formalise CML-driven rainfall nowcasting as a two-stage
        spatiotemporal learning problem -- a recurrent CML-to-rain
        estimator followed by a learned dynamical model on the resulting
        rain field -- and clarify the resulting upper-bound and bias
        properties.
  \item We adapt and benchmark three families of dynamical models
        (Transformer, POD--SINDy, recurrent SSM) within this framework
        on the public OpenMRG dataset~\cite{andersson2022openmrg}, using
        identical inputs, lookback, and supervision protocols.
  \item We study two orthogonal design axes that are usually conflated
        in the nowcasting literature: \emph{(i)} single-step (one
        specialised model per horizon) versus multi-horizon
        (one model emitting all horizons) prediction; and \emph{(ii)}
        dense self-supervision on the estimated rain field versus sparse
        gauge-anchored supervision at the loss level.
  \item We report a dual evaluation -- point-to-pixel against rain
        gauges and grid-to-grid against weather radar -- with
        continuous and detection-oriented metrics across four horizons
        (15--60\,min), and discuss implications for ISAC-driven
        environmental sensing.
\end{itemize}

% =============================================================================
\section{Related Work}
\label{sec:related}

\paragraph*{CML-based rain estimation}
The use of CML attenuation for rain monitoring dates back to
Messer~\emph{et~al.}~\cite{messer2006} and has since been refined with
careful baseline (dry-period) identification, wet-antenna attenuation
correction, and rain-rate retrieval~\cite{overeem2013,chwala2019}.
Recent learning-based estimators replace handcrafted pipelines with
recurrent or attention networks trained against gauge
references~\cite{habi2020,pynncml}; the \texttt{pynncml} toolkit and the
publicly released OpenMRG dataset~\cite{andersson2022openmrg} have made
reproducible benchmarking practical.

\paragraph*{Rainfall nowcasting}
Classical radar-based nowcasting techniques such as optical-flow
extrapolation and the Lagrangian persistence implemented in
PySTEPS~\cite{pulkkinen2019pysteps} remain strong baselines.
Deep-learning approaches -- ConvLSTM~\cite{shi2015convlstm},
PredRNN~\cite{wang2017predrnn}, and more recently NowcastNet and related
generative models -- have shown improvements on radar-only benchmarks.
With rare exceptions, none of these methods consume CML data, and the
spatial inductive biases they assume (dense, gridded radar input) are
not directly suitable for sparse, line-integrated CML observations.

\paragraph*{Dynamical sequence models}
Beyond rainfall, the choice between attention~\cite{vaswani2017},
selective state-space recurrence~\cite{gu2023mamba}, and
equation-discovery~\cite{brunton2016sindy,kaiser2018sindyc} is now a
recurring question in any spatiotemporal forecasting application.
POD--SINDy in particular has been used to identify reduced-order models
for fluid and atmospheric fields~\cite{loiseau2018sindyfluid}, but its
behaviour on CML-derived precipitation maps -- which are noisy,
sparsely informed, and intermittently zero -- has not been characterised.

% =============================================================================
\section{Problem Formulation}
\label{sec:problem}

Let $\mathcal{L}=\{1,\dots,L\}$ index a set of CMLs, each with a
recorded RSL/TSL time series $\mathbf{s}_\ell(\tau)\in\mathbb{R}^{2}$
and static metadata $\mathbf{m}_\ell\in\mathbb{R}^{d_m}$ (length,
frequency, polarisation, endpoint coordinates). Let
$\mathcal{G}=\{1,\dots,G\}$ index $G$ rain gauges providing point rain
rates $r^{\mathrm{g}}_g(\tau)$, and let $R^{\mathrm{rad}}(t,x,y)$ denote
the gridded weather-radar rain field on a regular $H\times W$ lattice.

\subsection{Stage~1: CML-to-Rain Estimation}

A learned mapping
$f_{\theta}:\mathbf{s}_\ell,\mathbf{m}_\ell \mapsto
\widehat{r}_\ell(\tau)$ produces a per-link rain estimate, trained
against gauge-derived link supervision in the standard manner. A fixed
inverse-distance-weighted (IDW) operator $\Phi$ then projects the
per-link estimates onto the radar grid, producing the CML-derived
rain map
\begin{equation}
  \widehat{R}^{\mathrm{cml}}(t,x,y) \;=\;
  \Phi\!\bigl(\{\widehat{r}_\ell(t)\}_{\ell\in\mathcal{L}}\bigr)(x,y).
\end{equation}
This map is the input substrate of the downstream forecaster; the radar
field is never used as a training target at this stage.

\subsection{Stage~2: Learned Spatiotemporal Dynamics}

Let $T_{\mathrm{lb}}$ denote the lookback length (we use
$T_{\mathrm{lb}}=8$ frames of 15\,min, i.e.~two hours of context).
Stage~2 learns a parametric dynamical operator
\begin{equation}
  g_{\phi}\colon\;
  \widehat{R}^{\mathrm{cml}}_{[t-T_{\mathrm{lb}}+1:t]}
  \;\longmapsto\;
  \widehat{R}^{\mathrm{pred}}_{t+h},
  \quad h\in\mathcal{H},
\label{eq:stage2}
\end{equation}
where $\mathcal{H}=\{1,2,3,4\}$ (i.e.~15, 30, 45, 60\,min ahead).
Equation~\eqref{eq:stage2} admits two natural realisations:
\emph{(a)} a separate operator $g_\phi^{(h)}$ per horizon (single-step
heads, $|\mathcal{H}|$ checkpoints), or \emph{(b)} a single operator
$g_\phi$ with a $|\mathcal{H}|\!\times\!H\!\times\!W$ output head
(multi-horizon).

\subsection{Supervision Variants}

Two loss specifications are considered. The \emph{self-supervised}
variant uses dense mean-squared error against the future CML-derived
field,
\begin{equation}
  \mathcal{L}_{\mathrm{self}} \;=\;
  \tfrac{1}{HW}\sum_{x,y}\bigl(
  \widehat{R}^{\mathrm{pred}}_{t+h}(x,y) - \widehat{R}^{\mathrm{cml}}_{t+h}(x,y)
  \bigr)^{2}.
\end{equation}
The \emph{gauge-supervised} variant computes the same map output but
restricts the loss to the $G$ gauge pixels with measured rain rate as
target,
\begin{equation}
  \mathcal{L}_{\mathrm{gauge}} \;=\;
  \tfrac{1}{G}\sum_{g}\bigl(
  \widehat{R}^{\mathrm{pred}}_{t+h}(x_g,y_g) - r^{\mathrm{g}}_g(t+h)
  \bigr)^{2}.
\end{equation}
The self-supervised variant offers dense gradient signal that is biased
toward whatever errors Stage~1 already commits; the gauge-supervised
variant grounds the model in physical measurements but provides only
$G\!\sim\!\mathcal{O}(10)$ supervision points per timestep.

% =============================================================================
\section{Dynamical Models}
\label{sec:methods}

We instantiate the operator $g_\phi$ in
Equation~\eqref{eq:stage2} with three architectures spanning the
attention, state-space, and equation-discovery axes. All three consume
the same lookback tensor of shape
$(T_{\mathrm{lb}},H,W)$ and emit either a single map or a stack of
$|\mathcal{H}|$ maps.

\subsection{Attention Dynamics: Map-Token Transformer}

We flatten each frame to a vector of dimension $HW$ and embed it
linearly to a $d_\mathrm{model}$-dimensional token, producing a sequence
of $T_{\mathrm{lb}}$ tokens. A standard pre-LN Transformer encoder with
$n_h$ heads and $n_\ell$ layers~\cite{vaswani2017} processes the
sequence, and a linear decoder maps the final-position embedding to one
or $|\mathcal{H}|$ flattened $H\!\times\!W$ outputs.

\subsection{Reduced-Order Equation Discovery: POD--SINDy}

We compute a truncated SVD of the training rain fields with $r=10$
modes, yielding a latent state
$\mathbf{z}(t)=\mathbf{U}^{\!\top}\mathrm{vec}\bigl(\widehat{R}^{\mathrm{cml}}(t)\bigr)$.
SINDy~\cite{brunton2016sindy} with a polynomial library of degree two
and a sequentially-thresholded least-squares (STLSQ) sparsity prior
identifies the dynamics
\begin{equation}
  \dot{\mathbf{z}}(t) \;=\; \mathbf{\Xi}\,\boldsymbol{\Theta}(\mathbf{z}(t)),
\end{equation}
where $\boldsymbol{\Theta}$ is the polynomial feature library and
$\mathbf{\Xi}$ is sparse. Forecasting reduces to integrating the
identified ODE for $h\!\cdot\!\Delta t$ from the initial condition
provided by the last lookback frame, followed by an inverse SVD lift to
the spatial grid. POD--SINDy is \emph{horizon-agnostic}: the same fitted
$\mathbf{\Xi}$ produces every horizon by integrating the trajectory to
the corresponding time. The gauge-supervised variant of POD--SINDy fits
the same pipeline to the IDW gauge field instead of
$\widehat{R}^{\mathrm{cml}}$.

\subsection{State-Space Dynamics: Recurrent SSM Surrogate}

In the spirit of selective state-space sequence
models~\cite{gu2023mamba}, we instantiate a recurrent surrogate that
processes the same flattened map-token sequence with a multi-layer gated
recurrent unit (GRU)~\cite{cho2014gru} and projects the final hidden
state to one or $|\mathcal{H}|$ output maps. We adopt the GRU formulation
as a tractable SSM-style recurrence whose hidden state plays the role of
the latent system state; replacing the GRU with the full
Mamba block is a straightforward drop-in and is left for the camera-ready
revision (see Section~\ref{sec:discussion}).

\subsection{Single-Step versus Multi-Horizon Heads}

For each architecture, the single-step variant trains $|\mathcal{H}|$
separate parameter sets, one per horizon, each minimising the loss for
its target horizon alone. The multi-horizon variant widens the output
head to emit all $|\mathcal{H}|$ horizons simultaneously and trains
under their averaged loss. The trade-off is classical: specialisation
versus parameter sharing across horizons.

% =============================================================================
\section{Experimental Setup}
\label{sec:setup}

\subsection{Dataset and Splits}

We use the public OpenMRG
dataset~\cite{andersson2022openmrg}, covering Gothenburg, Sweden, over
1~June -- 31~August 2015, comprising several hundred CMLs, a sparse rain
gauge network, and a co-located C-band radar product. All channels are
resampled to a common 15-minute cadence and cropped to a
lat~$[57.55, 57.85]$, lon~$[11.70, 12.20]$ window covering the
gauge-rich region. Radar reflectivity is converted to rain rate via the
Marshall--Palmer relation
$R=(Z/200)^{1/1.6}$~\cite{marshall1948}. Chronological splits are
Train (1~Jun~--~20~Jul), Val (21~Jul~--~5~Aug), and Test
(6~Aug~--~31~Aug); no shuffled split is used.

\subsection{Stage~1 Estimator}

The Stage~1 estimator is the two-step GRU rain estimator distributed
with \texttt{pynncml}~\cite{pynncml}, trained with a regression head for
rain rate and a binary head for wet/dry detection under the loss
weighting recommended therein. Hyperparameters are selected by a
ten-trial Bayesian sweep on (lr, batch, hidden size, depth, weight
decay), and the final estimator is trained for 200 epochs on the Train
split with selection by Val loss. The fitted estimator is then frozen
and run over the entire campaign to produce
$\widehat{R}^{\mathrm{cml}}$, which Stage~2 consumes.

\subsection{Stage~2 Hyperparameter Search}

For both the Transformer and the SSM surrogate, we run a six-trial
Bayesian sweep over architectural hyperparameters
($d_{\mathrm{model}}$, $n_h$, $n_\ell$ for the Transformer; hidden size
and depth for the SSM) and optimisation hyperparameters (learning rate,
batch size). Sweep trials train on horizon $h=1$ for 30 epochs;
best-by-Val-loss configurations are then retrained from scratch for
80~epochs at every horizon, with early stopping on Val loss.

\subsection{Baselines}

We compare against three reference predictors: \emph{(i)} naive radar
persistence, $\widehat{R}(t+h)=R^{\mathrm{rad}}(t-h)$;
\emph{(ii)} naive gauge persistence via IDW interpolation of past gauge
readings; and \emph{(iii)} the Stage~1 CML estimator $\widehat{R}^{\mathrm{cml}}(t)$
itself, treated as an $h=0$ no-forecast lower bound that any useful
predictor at $h>0$ should at least approach.

\subsection{Metrics}

We report continuous metrics (RMSE, MAE, mean bias, Pearson CC) in
mm/h and detection metrics at a wet/dry threshold of 0.1\,mm/h
(accuracy, F1, probability of detection POD, false-alarm ratio FAR,
critical success index CSI, and Heidke skill score HSS). Continuous
metrics summarise quantitative agreement; detection metrics summarise
event-level skill and HSS is particularly informative since it corrects
for the climatological wet/dry base rate.

Evaluation is performed twice. \emph{Point-to-pixel} compares the
predicted map, sampled at the nearest grid pixel to each gauge, against
the gauge reading -- this is the most direct physical evaluation but
involves only $G$ test points per frame. \emph{Grid-to-grid} compares
the full predicted map (cropped to the lat/lon window) against the
radar field, providing spatial coverage at the cost of using a
secondary reference.

% =============================================================================
\section{Results}
\label{sec:results}

\subsection{Headline Comparison Across Horizons}

Table~\ref{tab:headline-gauges} summarises the gauge-anchored
point-to-pixel performance of every model and baseline at all four
horizons. Table~\ref{tab:headline-radar} reports the corresponding
grid-to-grid evaluation against radar.

\begin{table}[t]
\centering
\caption{Point-to-pixel performance against rain gauges (test split).
Lower is better for RMSE, MAE, FAR; higher is better for the others.}
\label{tab:headline-gauges}
\begin{tabular}{l c c c c c c}
\toprule
Model & Horizon & RMSE & MAE & CC & F1 & HSS \\
      & (min)   & (mm/h) & (mm/h) & & & \\
\midrule
Naive radar    & 15 & \tbd & \tbd & \tbd & \tbd & \tbd \\
Naive gauge    & 15 & \tbd & \tbd & \tbd & \tbd & \tbd \\
Transformer    & 15 & \tbd & \tbd & \tbd & \tbd & \tbd \\
POD--SINDy     & 15 & \tbd & \tbd & \tbd & \tbd & \tbd \\
SSM surrogate  & 15 & \tbd & \tbd & \tbd & \tbd & \tbd \\
\midrule
Naive radar    & 30 & \tbd & \tbd & \tbd & \tbd & \tbd \\
Transformer    & 30 & \tbd & \tbd & \tbd & \tbd & \tbd \\
POD--SINDy     & 30 & \tbd & \tbd & \tbd & \tbd & \tbd \\
SSM surrogate  & 30 & \tbd & \tbd & \tbd & \tbd & \tbd \\
\midrule
\multicolumn{7}{c}{\emph{(15-min and 60-min rows analogous)}}\\
\bottomrule
\end{tabular}
\end{table}

\begin{table}[t]
\centering
\caption{Grid-to-grid performance against weather radar (test split).}
\label{tab:headline-radar}
\begin{tabular}{l c c c c c c}
\toprule
Model & Horizon & RMSE & MAE & CC & F1 & HSS \\
      & (min)   & (mm/h) & (mm/h) & & & \\
\midrule
\multicolumn{7}{c}{\emph{(populated identically to Tab.~\ref{tab:headline-gauges})}}\\
\bottomrule
\end{tabular}
\end{table}

\noindent\tbd{} \note{Narrative: pick the best learned model on the
gauge metric, quantify its margin over the strongest naive baseline at
each horizon, and identify the horizon at which learned models cross or
fail to cross the radar persistence baseline. Discuss whether the
ranking is consistent between gauge and radar evaluation.}

\subsection{Single-Step versus Multi-Horizon Heads}

Within each architecture family we contrast the single-step and
multi-horizon variants. \tbd{} \note{Narrative: report the
per-horizon RMSE delta (multi $-$ single) for each family; identify
whether one family benefits more than another from parameter sharing
across horizons; comment on training cost (single requires
$|\mathcal{H}|$ checkpoints).}

\subsection{Self-Supervision versus Gauge-Supervision}

Comparing the dense-MSE training against $\widehat{R}^{\mathrm{cml}}$
with the gauge-anchored loss isolates the effect of supervision
sparsity. \tbd{} \note{Narrative: hypothesis is that gauge-supervised
variants win the gauge evaluation but lose the radar evaluation; confirm
or refute. Comment on whether sparsity of supervision induces
visibly biased spatial maps (over-/under-prediction far from gauges).}

\subsection{Qualitative Inspection}

For the wettest test frame, all forecasters at the 30-minute horizon
produce maps that can be inspected against the radar truth.
\tbd{} \note{Describe in words which model best localises the storm,
which over- or under-predicts intensity, and whether POD--SINDy's
low-rank reconstruction visibly smears spatial structure.}

% =============================================================================
\section{Discussion}
\label{sec:discussion}

\paragraph*{Stage~1 as a performance ceiling}
A key consequence of the two-stage formulation is that the Stage~2
forecasters can never outperform the underlying Stage~1 estimator at
$h\!=\!0$. In particular, the dense self-supervised variant explicitly
treats $\widehat{R}^{\mathrm{cml}}$ as ground truth, propagating any
Stage~1 bias into the forecast. Reporting the Stage~1 error against
gauges and radar (as we do) makes this ceiling explicit and clarifies
how much of any residual nowcasting error is attributable to estimation
versus to dynamical prediction.

\paragraph*{Sparsity-induced trade-offs}
Gauge supervision grounds the loss in direct measurements but provides
only $G\!\sim\!\mathcal{O}(10)$ supervision points per timestep, so the
trained map is only weakly constrained away from gauges. Conversely,
self-supervision provides dense gradient signal at every pixel but
optimises against an imperfect proxy. Our two-axis ablation provides
the practitioner with empirical guidance on which trade-off is
preferable in a given deployment.

\paragraph*{Inductive biases of the three families}
The Transformer imposes the weakest prior and is most data-hungry;
POD--SINDy imposes the strongest prior (low-rank latent state with
sparse polynomial dynamics) and is the most interpretable; the recurrent
SSM sits between the two. We expect their relative ranking to depend on
training-set size, on the variability of weather regimes seen at
training time, and on the dominant spatial scale of the storms in the
test split.

\paragraph*{Limitations}
The present study is restricted to one campaign (Gothenburg summer 2015)
with a single radar reference and roughly ten gauges; the variance of
gauge-anchored metrics is therefore non-negligible and we report
bootstrap confidence intervals where appropriate. The recurrent SSM
surrogate uses a GRU rather than the full Mamba block; replacing the
GRU with a Mamba layer is a drop-in change and we expect it to improve
long-context behaviour. Finally, an end-to-end training of Stage~1 and
Stage~2 jointly -- bypassing the intermediate
$\widehat{R}^{\mathrm{cml}}$ representation -- is an obvious next step,
at the cost of architectural complexity and reduced interpretability.

% =============================================================================
\section{Conclusion}
\label{sec:conclusion}

We have presented a two-stage CML-driven rainfall nowcasting framework
and benchmarked three families of data-driven dynamical models -- a
Transformer, a POD--SINDy reduced-order model, and a recurrent
state-space surrogate -- against persistence baselines on the OpenMRG
dataset. Our results indicate that \tbd{}. More broadly, the work
illustrates that opportunistic CML sensing, combined with modern
sequence models, is a viable ISAC pathway toward short-term rainfall
nowcasting at deployment scales where dedicated weather radar is not
economic. Future work will incorporate the full Mamba block, larger
multi-season campaigns, ensemble-based uncertainty quantification, and
end-to-end joint training of estimation and forecasting.

% =============================================================================
\section*{Acknowledgments}
\note{Funding, dataset access, and computing acknowledgments here.}

% =============================================================================
\bibliographystyle{IEEEtran}
\bibliography{bibs/refs, bibs/essd}

\end{document}
